\documentclass[11pt]{article}
\usepackage[margin=1in]{geometry}
\usepackage{enumitem}
\usepackage{array}
\usepackage{booktabs}
\usepackage{xcolor}
\usepackage{tcolorbox}

\title{\textbf{Trauma, Mass Casualty \& Spinal Cord Injuries}\\Fill-in-the-Blank Study Activity}
\author{Med Surg 2 Exam Review}
\date{}

\begin{document}
\maketitle

\section*{Instructions}
Fill in each blank without looking at your notes. After completing all sections, check your answers using the key at the end. Mark any items you missed for focused review.

\hrulefill

\section{Trauma Overview}

\subsection{Definitions}
\begin{enumerate}
    \item Trauma is the leading cause of death for Americans aged \underline{\hspace{2cm}}.
    
    \item \textbf{Blunt trauma:} No break in skin (examples: MVA, falls, assaults); injury from compression, \underline{\hspace{2cm}}, and deceleration forces.
    
    \item \textbf{Penetrating trauma:} Break in skin (examples: \underline{\hspace{3cm}}, stabbings); injury depends on trajectory and \underline{\hspace{2cm}}.
\end{enumerate}

\subsection{Primary Survey - ABCDE}
Write what each letter stands for:

\begin{enumerate}
    \item \textbf{A} = \underline{\hspace{4cm}} with cervical spine protection
    \item \textbf{B} = \underline{\hspace{4cm}} and ventilation
    \item \textbf{C} = \underline{\hspace{4cm}} with hemorrhage control
    \item \textbf{D} = \underline{\hspace{4cm}} (neurological status - GCS, pupils)
    \item \textbf{E} = \underline{\hspace{4cm}} / Environmental control
\end{enumerate}

\subsection{Secondary Survey}
The secondary survey includes head-to-toe assessment and history using the mnemonic \textbf{AMPLE}. What does each letter represent?

\begin{itemize}
    \item A = \underline{\hspace{3cm}}
    \item M = \underline{\hspace{3cm}}
    \item P = \underline{\hspace{3cm}}
    \item L = \underline{\hspace{3cm}}
    \item E = \underline{\hspace{3cm}}
\end{itemize}

\hrulefill

\section{Hemorrhagic Shock Classes}

Complete the table:

\begin{center}
\begin{tabular}{|c|c|c|c|c|}
\hline
\textbf{Class} & \textbf{Blood Loss} & \textbf{HR} & \textbf{BP} & \textbf{Mental Status} \\
\hline
Class I & Up to 750 mL ($<$15\%) & $<$100 & \underline{\hspace{1.5cm}} & Normal \\
\hline
Class II & \underline{\hspace{2.5cm}} & 100-120 & Normal & \underline{\hspace{2cm}} \\
\hline
Class III & 1500-2000 mL (30-40\%) & \underline{\hspace{1.5cm}} & Decreased & Confused \\
\hline
Class IV & $>$2000 mL ($>$40\%) & $>$140 & \underline{\hspace{1.5cm}} & \underline{\hspace{2.5cm}} \\
\hline
\end{tabular}
\end{center}

\hrulefill

\section{Trauma Nursing Priorities}

\begin{enumerate}
    \item Establish large-bore IV access ($\times$2, \underline{\hspace{2cm}}-gauge or larger)
    
    \item Warm crystalloid fluid resuscitation (\underline{\hspace{2cm}} preferred; avoid hypothermia)
    
    \item Massive transfusion protocol (MTP) ratio for packed RBCs : FFP : platelets = \underline{\hspace{2cm}}
    
    \item \textbf{Permissive hypotension} for penetrating trauma: target SBP \underline{\hspace{2cm}}
    
    \item The ``trauma triad of death'' includes: hypothermia, \underline{\hspace{2cm}}, and coagulopathy
\end{enumerate}

\hrulefill

\section{Life-Threatening Thoracic Injuries (``Deadly Dozen'')}

Match the injury with its key features/treatment:

\begin{enumerate}
    \item \textbf{Tension pneumothorax:}
    \begin{itemize}
        \item Trachea deviated \underline{\hspace{2cm}} from injury
        \item Treatment: IMMEDIATE needle decompression at \underline{\hspace{2cm}} ICS, MCL
    \end{itemize}
    
    \item \textbf{Open pneumothorax (sucking chest wound):}
    \begin{itemize}
        \item Treatment: \underline{\hspace{2cm}}-sided dressing, then chest tube
    \end{itemize}
    
    \item \textbf{Massive hemothorax:}
    \begin{itemize}
        \item Defined as $>$ \underline{\hspace{2cm}} mL blood in chest
    \end{itemize}
    
    \item \textbf{Cardiac tamponade (Beck's triad):}
    \begin{itemize}
        \item JVD + \underline{\hspace{2cm}} + \underline{\hspace{3cm}}
        \item Treatment: \underline{\hspace{3cm}}
    \end{itemize}
    
    \item \textbf{Flail chest:}
    \begin{itemize}
        \item Key sign: \underline{\hspace{3cm}} chest wall movement
    \end{itemize}
\end{enumerate}

\hrulefill

\section{Mass Casualty Incidents (MCI)}

\subsection{START Triage System}
START = \underline{\hspace{6cm}}

Time limit per patient: \underline{\hspace{2cm}} seconds

Complete the triage steps:

\begin{enumerate}
    \item Can the patient walk? $\rightarrow$ Yes = \underline{\hspace{2cm}} (Minor/Delayed)
    
    \item Is the patient breathing? $\rightarrow$ No, reposition airway $\rightarrow$ Still no = \underline{\hspace{2cm}} (Expectant/Deceased)
    
    \item Respiratory rate $>$30 or $<$10 = \underline{\hspace{2cm}} (Immediate)
    
    \item Radial pulse absent OR capillary refill $>$ \underline{\hspace{1cm}} sec = RED
    
    \item Cannot follow commands = \underline{\hspace{2cm}}
    
    \item If steps 3-5 are normal = \underline{\hspace{2cm}} (Delayed)
\end{enumerate}

\subsection{Triage Colors}
\begin{center}
\begin{tabular}{|c|c|p{7cm}|}
\hline
\textbf{Color} & \textbf{Priority} & \textbf{Meaning} \\
\hline
RED & \underline{\hspace{2cm}} & Life-threatening; can survive with immediate intervention \\
\hline
YELLOW & Delayed & \underline{\hspace{5cm}} \\
\hline
GREEN & \underline{\hspace{2cm}} & Walking wounded; treatment can be delayed \\
\hline
BLACK & Expectant & \underline{\hspace{5cm}} \\
\hline
\end{tabular}
\end{center}

\hrulefill

\section{Spinal Cord Injuries (SCI)}

\subsection{Definitions}
\begin{enumerate}
    \item \textbf{Complete SCI:} \underline{\hspace{4cm}} loss of motor and sensory function below injury level
    
    \item \textbf{Incomplete SCI:} \underline{\hspace{4cm}} function preserved below injury level
    
    \item Most common cause of SCI: \underline{\hspace{3cm}}
\end{enumerate}

\subsection{Levels of SCI and Functional Outcomes}
Complete the table:

\begin{center}
\begin{tabular}{|c|c|p{6cm}|}
\hline
\textbf{Level} & \textbf{Classification} & \textbf{Expected Function} \\
\hline
C1-C4 & High cervical & \underline{\hspace{4cm}}; minimal to no limb function \\
\hline
C5-C8 & Lower cervical & \underline{\hspace{3cm}}; partial arm use possible \\
\hline
T1-T12 & Thoracic & \underline{\hspace{3cm}}; full arm use; respiratory may be affected \\
\hline
L1-L5 & Lumbar & Paraplegia; \underline{\hspace{3cm}} \\
\hline
S1-S5 & Sacral & \underline{\hspace{4cm}}; leg weakness \\
\hline
\end{tabular}
\end{center}

\subsection{Memory Device}
``C3, 4, 5 keeps the \underline{\hspace{3cm}} alive'' - these levels innervate the diaphragm

\hrulefill

\section{Spinal Shock vs. Neurogenic Shock}

\begin{center}
\begin{tabular}{|p{3cm}|p{4cm}|p{4cm}|}
\hline
\textbf{Feature} & \textbf{Spinal Shock} & \textbf{Neurogenic Shock} \\
\hline
Definition & Temporary suppression of \underline{\hspace{2cm}} & Hemodynamic instability from \underline{\hspace{2cm}} loss \\
\hline
Duration & \underline{\hspace{2cm}} & Ongoing until spinal stability \\
\hline
Hemodynamics & Normal & Hypotension + \underline{\hspace{2cm}} (no compensatory tachycardia) \\
\hline
Temperature & Normal & \underline{\hspace{2cm}} (cannot vasoconstrict) \\
\hline
Treatment & Supportive & \underline{\hspace{2cm}}, atropine, warming \\
\hline
\end{tabular}
\end{center}

\hrulefill

\section{Autonomic Dysreflexia}

\textbf{CRITICAL:} Medical emergency in SCI at or above \underline{\hspace{2cm}}

\subsection{Triggers (most common first)}
\begin{enumerate}
    \item \underline{\hspace{3cm}} (most common)
    \item Bowel impaction
    \item Pressure injury
    \item Tight clothing
\end{enumerate}

\subsection{Signs/Symptoms}
\begin{itemize}
    \item Sudden severe hypertension (SBP $>$ \underline{\hspace{2cm}} mmHg)
    \item Pounding \underline{\hspace{2cm}}
    \item Flushing/sweating \underline{\hspace{2cm}} injury level
    \item Pallor \underline{\hspace{2cm}} injury level
    \item \underline{\hspace{2cm}} (NOT tachycardia!)
    \item Nasal congestion
\end{itemize}

\subsection{Immediate Actions (in order)}
\begin{enumerate}
    \item Sit patient \underline{\hspace{2cm}} (lowers BP by orthostatic effect)
    \item Loosen all \underline{\hspace{3cm}}
    \item Check and empty \underline{\hspace{2cm}} (kink/obstruction in catheter?)
    \item Check rectum for \underline{\hspace{2cm}} (use anesthetic gel first)
    \item Administer antihypertensive: \underline{\hspace{2cm}}
    \item Monitor every \underline{\hspace{2cm}} minutes; notify provider
\end{enumerate}

\hrulefill

\section{Acute SCI Nursing Priorities}

\begin{enumerate}
    \item \textbf{Immobilization:} Rigid cervical collar + \underline{\hspace{2cm}} until injury ruled out; \underline{\hspace{2cm}} technique for repositioning
    
    \item \textbf{Respiratory:} Monitor closely for cervical SCI; prepare for \underline{\hspace{2cm}}
    
    \item \textbf{Neurogenic shock treatment:} IV fluids, \underline{\hspace{2cm}}, atropine for bradycardia
    
    \item \textbf{Prevent secondary injury:} Maintain MAP $>$ \underline{\hspace{2cm}}; avoid hypoxia, hyperthermia
    
    \item \textbf{Bladder:} \underline{\hspace{3cm}} catheterization; monitor for autonomic dysreflexia
    
    \item \textbf{Bowel:} Bowel program; \underline{\hspace{2cm}}; prevent impaction
    
    \item \textbf{Skin:} Reposition every \underline{\hspace{2cm}} hours
    
    \item \textbf{DVT prophylaxis:} SCDs, \underline{\hspace{3cm}} heparin
\end{enumerate}

\hrulefill

\section{NGN Practice: Primary Survey Priority Order}

Place these interventions in correct priority order (1-7):

\begin{tabular}{|p{0.5cm}|p{12cm}|}
\hline
\underline{\hspace{0.5cm}} & Apply manual in-line cervical spine stabilization and assess airway patency \\
\hline
\underline{\hspace{0.5cm}} & Perform needle decompression of the left chest (2nd ICS, MCL) \\
\hline
\underline{\hspace{0.5cm}} & Apply direct pressure to forehead laceration \\
\hline
\underline{\hspace{0.5cm}} & Establish two large-bore IVs and initiate LR fluid resuscitation \\
\hline
\underline{\hspace{0.5cm}} & Assess GCS, pupils, and motor strength bilaterally \\
\hline
\underline{\hspace{0.5cm}} & Expose patient fully and apply warm blankets to prevent hypothermia \\
\hline
\underline{\hspace{0.5cm}} & Assess for bilateral breath sounds and chest excursion \\
\hline
\end{tabular}

\vspace{0.5cm}
\textbf{Rationale:} Why does needle decompression come BEFORE IV access?

\underline{\hspace{12cm}}

\underline{\hspace{12cm}}

\newpage
\section*{ANSWER KEY}

\subsection*{Section 1: Trauma Overview}
\begin{enumerate}
    \item 1-44
    \item shear
    \item gunshot wounds; energy
\end{enumerate}

\textbf{Primary Survey:}
A = Airway; B = Breathing; C = Circulation; D = Disability; E = Exposure

\textbf{AMPLE:}
A = Allergies; M = Medications; P = Past medical history; L = Last meal; E = Events

\subsection*{Section 2: Hemorrhagic Shock}
Class I: BP = Normal\\
Class II: 750-1500 mL (15-30\%); Anxious\\
Class III: HR = 120-140\\
Class IV: BP = Very low; Lethargic/unconscious

\subsection*{Section 3: Trauma Nursing Priorities}
\begin{enumerate}
    \item 16-gauge
    \item LR (Lactated Ringer's)
    \item 1:1:1
    \item 80-90
    \item acidosis
\end{enumerate}

\subsection*{Section 4: Thoracic Injuries}
\begin{enumerate}
    \item AWAY; 2nd
    \item Three
    \item 1500
    \item Hypotension; Muffled heart sounds; Pericardiocentesis
    \item Paradoxical
\end{enumerate}

\subsection*{Section 5: MCI/START Triage}
START = Simple Triage and Rapid Treatment; 60 seconds

Steps: 1. GREEN; 2. BLACK; 3. RED; 4. 2 sec; 5. RED; 6. YELLOW

Colors: RED = Immediate; YELLOW = Serious but stable, can wait 30-60 min; GREEN = Minor; BLACK = Deceased or unsurvivable

\subsection*{Section 6: SCI}
\begin{enumerate}
    \item Total
    \item Some; better
    \item MVA
\end{enumerate}

\subsection*{Section 7: SCI Levels}
C1-C4 = Ventilator-dependent; minimal to no limb function\\
C5-C8 = Quadriplegia; partial arm use (C6-C8)\\
T1-T12 = Paraplegia; full arm use; respiratory may be affected\\
L1-L5 = Paraplegia; variable leg function\\
S1-S5 = Bowel/bladder dysfunction; leg weakness

\subsection*{Section 8: Spinal Shock vs. Neurogenic Shock}
\textbf{Spinal Shock:}
Definition = Temporary suppression of all reflexes\\
Duration = Hours to weeks\\
Hemodynamics = Normal\\
Temperature = Normal\\
Treatment = Supportive

\textbf{Neurogenic Shock:}
Definition = Hemodynamic instability from sympathetic loss\\
Duration = Ongoing until spinal stability\\
Hemodynamics = Hypotension + bradycardia (no compensatory tachycardia)\\
Temperature = Hypothermia (cannot vasoconstrict)\\
Treatment = Vasopressors (norepinephrine), atropine, warming

\subsection*{Section 9: Autonomic Dysreflexia}
Occurs at or above = T6\\
Triggers = Bladder distension (most common), bowel impaction, pressure injury, tight clothing\\
Signs = SBP > 150 mmHg, pounding headache, flushing/sweating ABOVE injury, pallor BELOW injury, bradycardia, nasal congestion

\textbf{Immediate Actions:}
\begin{enumerate}
    \item Sit patient upright (orthostatic effect lowers BP)
    \item Loosen restrictive clothing
    \item Check and empty bladder
    \item Check rectum for impaction (use anesthetic gel first)
    \item Administer nifedipine if BP remains elevated
    \item Monitor every 5 minutes; notify provider
\end{enumerate}

\subsection*{Section 10: Acute SCI Nursing Priorities}
\begin{enumerate}
    \item Immobilization: Rigid cervical collar + backboard; log-roll technique
    \item Respiratory: C3-C5 (``C3, 4, 5 keeps the diaphragm alive'')
    \item Neurogenic shock: IV fluids, vasopressors, atropine for bradycardia
    \item Prevent secondary injury: MAP > 85, avoid hypoxia, hyperthermia
    \item Bladder: Intermittent catheterization; monitor for autonomic dysreflexia
    \item Bowel: Bowel program; stool softeners; prevent impaction
    \item Skin: Reposition every 2 hours
    \item DVT prophylaxis: SCDs, low-molecular-weight heparin
    \item Psychological support: Grief counseling; address body image and life changes
\end{enumerate}

\subsection*{Section 11: NGN Priority Ordering}
Correct order: A $\rightarrow$ G $\rightarrow$ B $\rightarrow$ D $\rightarrow$ E $\rightarrow$ C $\rightarrow$ F

\begin{enumerate}
    \item A (Airway + C-spine) - always first in trauma
    \item G (Breathing assessment) - absent breath sounds = life threat
    \item B (Needle decompression) - tension pneumothorax is immediate life threat, treat BEFORE IV access
    \item D (Circulation) - IV access and fluid resuscitation for hemorrhagic shock
    \item E (Disability) - GCS and neurological assessment
    \item C (Hemorrhage control) - forehead laceration addressed after life threats
    \item F (Exposure/Environment) - prevent hypothermia
\end{enumerate}

\textbf{Key Rationale:} Tension pneumothorax (absent breath sounds + tracheal deviation + JVD + hemodynamic instability) is a Breathing emergency treated with needle decompression BEFORE any imaging or IV access.

\newpage
\section*{Section 12: SBAR Communication Practice}

\begin{tcolorbox}[colback=blue!5!white,colframe=blue!75!black,title=SBAR Scenario: Autonomic Dysreflexia]
Complete the SBAR communication for the following patient:

\textbf{Patient:} Mr. Garcia, 34-year-old in room 508 with T4 complete SCI, admitted 3 days ago after motorcycle accident. Foley catheter clamped 2 hours ago for bladder training. Now reporting sudden severe headache, BP 178/102 (baseline 100/70), flushed above nipple line, sweating profusely.
\end{tcolorbox}

\textbf{S (Situation):} Mr. Garcia, 34-year-old in room 508 with \underline{\hspace{2cm}} complete SCI, is reporting a sudden severe \underline{\hspace{2cm}} and his blood pressure is \underline{\hspace{2cm}}, up from his baseline of \underline{\hspace{2cm}}.

\vspace{0.3cm}

\textbf{B (Background):} He was admitted \underline{\hspace{2cm}} days ago after a \underline{\hspace{3cm}} accident. His Foley catheter was clamped \underline{\hspace{2cm}} hours ago for a \underline{\hspace{3cm}} trial. He also appears flushed above his \underline{\hspace{2cm}} and is sweating profusely.

\vspace{0.3cm}

\textbf{A (Assessment):} I believe this is \underline{\hspace{4cm}}, most likely triggered by \underline{\hspace{4cm}} from the clamped catheter.

\vspace{0.3cm}

\textbf{R (Recommendation):} I have already sat him \underline{\hspace{2cm}} and unclamped the \underline{\hspace{2cm}} -- BP is now coming down. I need an order for \underline{\hspace{3cm}} 10 mg sublingual in case it rises again. Please come evaluate.

\vspace{1cm}

\section*{Section 13: Critical Thinking - Clinical Judgment Scenarios}

\subsection*{Scenario A: Tension Pneumothorax Recognition}

\begin{tcolorbox}[colback=red!5!white,colframe=red!75!black,title=Case Presentation]
A 28-year-old male is brought to the ED after a high-speed motorcycle collision at approximately 60 mph. He was wearing a helmet. EMS reports he was found unconscious at the scene; he regained consciousness but is confused.

\textbf{On arrival:}
\begin{itemize}
    \item GCS 12 (Eyes 3, Verbal 3, Motor 6)
    \item BP 88/56 mmHg
    \item HR 128 bpm
    \item RR 30
    \item SpO2 90\% on 15L non-rebreather mask
\end{itemize}

\textbf{Assessment findings:}
\begin{itemize}
    \item Neck pain and bilateral leg numbness
    \item Trachea deviated to the RIGHT
    \item Absent breath sounds on the LEFT
    \item Distended neck veins (JVD)
    \item Rigid abdomen
    \item Pelvis unstable on compression
    \item 2.5 cm laceration on forehead
\end{itemize}
\end{tcolorbox}

\textbf{Question 1:} Based on the assessment findings, what is the MOST immediate life-threatening condition?

\underline{\hspace{12cm}}

\vspace{0.3cm}

\textbf{Question 2:} List the classic triad of findings that support your answer:

\begin{enumerate}[label=\alph*.]
    \item \underline{\hspace{8cm}}
    \item \underline{\hspace{8cm}}
    \item \underline{\hspace{8cm}}
\end{enumerate}

\vspace{0.3cm}

\textbf{Question 3:} What is the IMMEDIATE intervention required?

\underline{\hspace{12cm}}

\vspace{0.3cm}

\textbf{Question 4:} At what anatomical location is this procedure performed?

\underline{\hspace{6cm}} ICS, \underline{\hspace{6cm}} line

\vspace{0.3cm}

\textbf{Question 5:} Why is this intervention performed BEFORE establishing IV access?

\underline{\hspace{12cm}}

\underline{\hspace{12cm}}

\vspace{1cm}

\subsection*{Scenario B: Hemorrhagic Shock Classification}

\begin{tcolorbox}[colback=yellow!5!white,colframe=yellow!75!black,title=Patient Data]
Match each patient presentation to the correct hemorrhagic shock class:
\end{tcolorbox}

\textbf{Patient 1:} HR 95, BP 118/78, alert and oriented, estimated blood loss 600 mL

Class: \underline{\hspace{3cm}}

\vspace{0.3cm}

\textbf{Patient 2:} HR 135, BP 70/40, lethargic and barely responsive, massive bleeding from femoral artery

Class: \underline{\hspace{3cm}}

\vspace{0.3cm}

\textbf{Patient 3:} HR 115, BP 100/70, anxious and restless, estimated blood loss 1000 mL

Class: \underline{\hspace{3cm}}

\vspace{0.3cm}

\textbf{Patient 4:} HR 125, BP 85/60, confused, estimated blood loss 1800 mL

Class: \underline{\hspace{3cm}}

\vspace{1cm}

\section*{Section 14: START Triage Practice Scenarios}

\begin{tcolorbox}[colback=green!5!white,colframe=green!75!black,title=Mass Casualty Incident]
A tornado has struck a community center during an event. You are the triage nurse. Using the START system, assign the appropriate triage color to each victim. You have 60 seconds per patient.
\end{tcolorbox}

\textbf{Victim 1:} 45-year-old female, walking toward you asking for help, minor cuts on arms

\hspace{1cm} Can patient walk? \underline{\hspace{2cm}} $\rightarrow$ Triage Color: \underline{\hspace{3cm}}

\vspace{0.5cm}

\textbf{Victim 2:} 60-year-old male, not breathing, you reposition airway, still no breathing

\hspace{1cm} Breathing after repositioning? \underline{\hspace{2cm}} $\rightarrow$ Triage Color: \underline{\hspace{3cm}}

\vspace{0.5cm}

\textbf{Victim 3:} 30-year-old female, RR 35, radial pulse present, follows commands

\hspace{1cm} RR $>$ 30 or $<$ 10? \underline{\hspace{2cm}} $\rightarrow$ Triage Color: \underline{\hspace{3cm}}

\vspace{0.5cm}

\textbf{Victim 4:} 25-year-old male, RR 22, no radial pulse, capillary refill 4 seconds

\hspace{1cm} Radial pulse/cap refill $>$ 2 sec? \underline{\hspace{2cm}} $\rightarrow$ Triage Color: \underline{\hspace{3cm}}

\vspace{0.5cm}

\textbf{Victim 5:} 50-year-old female, RR 18, radial pulse present, does not follow commands, confused

\hspace{1cm} Mental status - follows commands? \underline{\hspace{2cm}} $\rightarrow$ Triage Color: \underline{\hspace{3cm}}

\vspace{0.5cm}

\textbf{Victim 6:} 8-year-old child, RR 24, radial pulse strong, follows commands, compound fracture of tibia

\hspace{1cm} All parameters normal? \underline{\hspace{2cm}} $\rightarrow$ Triage Color: \underline{\hspace{3cm}}

\vspace{1cm}

\section*{Section 15: SCI Level Matching Exercise}

\begin{tcolorbox}[colback=purple!5!white,colframe=purple!75!black,title=Instructions]
Match each patient presentation with the most likely level of spinal cord injury.
\end{tcolorbox}

\begin{tabular}{|p{8cm}|p{5cm}|}
\hline
\textbf{Patient Presentation} & \textbf{SCI Level} \\
\hline
Patient is ventilator-dependent with no limb movement & \underline{\hspace{4cm}} \\
\hline
Patient has full arm use but paraplegia; some respiratory compromise & \underline{\hspace{4cm}} \\
\hline
Patient has quadriplegia but can use wrists for some activities & \underline{\hspace{4cm}} \\
\hline
Patient has bowel/bladder dysfunction with leg weakness but can ambulate & \underline{\hspace{4cm}} \\
\hline
Patient has paraplegia with variable leg function, can transfer independently & \underline{\hspace{4cm}} \\
\hline
\end{tabular}

\vspace{0.5cm}

\textbf{Memory Aid:} Complete the phrase for cervical SCI respiratory function:

``C\underline{\hspace{0.5cm}}, \underline{\hspace{0.5cm}}, \underline{\hspace{0.5cm}} keeps the \underline{\hspace{3cm}} alive!''

\vspace{1cm}

\section*{Section 16: Differentiation Exercise - Spinal Shock vs. Neurogenic Shock}

For each characteristic, indicate whether it applies to \textbf{Spinal Shock (SS)}, \textbf{Neurogenic Shock (NS)}, or \textbf{Both (B)}:

\vspace{0.5cm}

\begin{tabular}{|p{10cm}|p{3cm}|}
\hline
\textbf{Characteristic} & \textbf{SS, NS, or B} \\
\hline
Occurs after spinal cord injury & \underline{\hspace{2.5cm}} \\
\hline
Temporary suppression of all reflexes below injury & \underline{\hspace{2.5cm}} \\
\hline
Hypotension with bradycardia (no compensatory tachycardia) & \underline{\hspace{2.5cm}} \\
\hline
Lasts hours to weeks & \underline{\hspace{2.5cm}} \\
\hline
Hemodynamics remain normal & \underline{\hspace{2.5cm}} \\
\hline
Requires vasopressors (norepinephrine) and atropine & \underline{\hspace{2.5cm}} \\
\hline
Patient is hypothermic (cannot vasoconstrict) & \underline{\hspace{2.5cm}} \\
\hline
Treatment is primarily supportive & \underline{\hspace{2.5cm}} \\
\hline
Results from loss of sympathetic tone & \underline{\hspace{2.5cm}} \\
\hline
Ongoing until spinal stability achieved & \underline{\hspace{2.5cm}} \\
\hline
\end{tabular}

\vspace{1cm}

\section*{Section 17: Autonomic Dysreflexia Action Sequence}

\begin{tcolorbox}[colback=red!10!white,colframe=red!75!black,title=Critical Alert]
Autonomic dysreflexia is a \textbf{MEDICAL EMERGENCY} in patients with SCI at or above T6. Failure to treat can result in stroke, seizures, or death.
\end{tcolorbox}

\textbf{Number the following nursing actions in the correct order (1-6):}

\vspace{0.3cm}

\begin{tabular}{|c|p{11cm}|}
\hline
\underline{\hspace{0.8cm}} & Administer antihypertensive if BP remains elevated (nifedipine) \\
\hline
\underline{\hspace{0.8cm}} & Check and empty bladder (check for kink/obstruction in catheter) \\
\hline
\underline{\hspace{0.8cm}} & Sit patient upright (lowers BP by orthostatic effect) \\
\hline
\underline{\hspace{0.8cm}} & Monitor every 5 minutes; notify provider \\
\hline
\underline{\hspace{0.8cm}} & Check rectum for impaction (use anesthetic gel first) \\
\hline
\underline{\hspace{0.8cm}} & Loosen all restrictive clothing/devices \\
\hline
\end{tabular}

\vspace{0.5cm}

\textbf{Critical Thinking:} Why must you use anesthetic gel BEFORE checking for fecal impaction?

\underline{\hspace{12cm}}

\underline{\hspace{12cm}}

\vspace{1cm}

\section*{Section 18: Comprehensive Review - True or False}

Mark each statement as \textbf{True (T)} or \textbf{False (F)}. If false, write the correct information.

\vspace{0.5cm}

\begin{enumerate}
    \item \underline{\hspace{1cm}} Trauma is the leading cause of death for Americans aged 45-64.
    
    Correction (if false): \underline{\hspace{8cm}}
    
    \vspace{0.3cm}
    
    \item \underline{\hspace{1cm}} In blunt trauma, there is a break in the skin from external force.
    
    Correction (if false): \underline{\hspace{8cm}}
    
    \vspace{0.3cm}
    
    \item \underline{\hspace{1cm}} The ``trauma triad of death'' includes hypothermia, acidosis, and coagulopathy.
    
    Correction (if false): \underline{\hspace{8cm}}
    
    \vspace{0.3cm}
    
    \item \underline{\hspace{1cm}} In tension pneumothorax, the trachea deviates TOWARD the affected side.
    
    Correction (if false): \underline{\hspace{8cm}}
    
    \vspace{0.3cm}
    
    \item \underline{\hspace{1cm}} Beck's triad for cardiac tamponade includes JVD, hypertension, and muffled heart sounds.
    
    Correction (if false): \underline{\hspace{8cm}}
    
    \vspace{0.3cm}
    
    \item \underline{\hspace{1cm}} During MCI triage, you should treat patients as you triage them.
    
    Correction (if false): \underline{\hspace{8cm}}
    
    \vspace{0.3cm}
    
    \item \underline{\hspace{1cm}} GREEN triage tag means the patient is deceased or has unsurvivable injuries.
    
    Correction (if false): \underline{\hspace{8cm}}
    
    \vspace{0.3cm}
    
    \item \underline{\hspace{1cm}} Permissive hypotension targets SBP of 80-90 mmHg in penetrating trauma.
    
    Correction (if false): \underline{\hspace{8cm}}
    
    \vspace{0.3cm}
    
    \item \underline{\hspace{1cm}} Autonomic dysreflexia occurs in SCI patients with injuries at or above T6.
    
    Correction (if false): \underline{\hspace{8cm}}
    
    \vspace{0.3cm}
    
    \item \underline{\hspace{1cm}} In neurogenic shock, the patient presents with tachycardia as a compensatory mechanism.
    
    Correction (if false): \underline{\hspace{8cm}}
    
\end{enumerate}

\vspace{1cm}

\section*{Section 19: Rapid-Fire Recall Questions}

Answer each question in the space provided:

\vspace{0.5cm}

\begin{enumerate}
    \item What does ABCDE stand for in the primary survey?
    
    A: \underline{\hspace{8cm}}
    
    B: \underline{\hspace{8cm}}
    
    C: \underline{\hspace{8cm}}
    
    D: \underline{\hspace{8cm}}
    
    E: \underline{\hspace{8cm}}
    
    \vspace{0.3cm}
    
    \item What does AMPLE stand for in the secondary survey history?
    
    A: \underline{\hspace{8cm}}
    
    M: \underline{\hspace{8cm}}
    
    P: \underline{\hspace{8cm}}
    
    L: \underline{\hspace{8cm}}
    
    E: \underline{\hspace{8cm}}
    
    \vspace{0.3cm}
    
    \item What does START stand for in triage?
    
    \underline{\hspace{10cm}}
    
    \vspace{0.3cm}
    
    \item What is the MTP ratio for packed RBCs : FFP : Platelets?
    
    \underline{\hspace{3cm}} : \underline{\hspace{3cm}} : \underline{\hspace{3cm}}
    
    \vspace{0.3cm}
    
    \item What size IV catheter is required for trauma resuscitation?
    
    \underline{\hspace{3cm}} gauge or larger, quantity: \underline{\hspace{2cm}}
    
    \vspace{0.3cm}
    
    \item What is the preferred crystalloid fluid for trauma resuscitation?
    
    \underline{\hspace{6cm}}
    
    \vspace{0.3cm}
    
    \item At what blood loss volume is massive hemothorax defined?
    
    $>$ \underline{\hspace{3cm}} mL
    
    \vspace{0.3cm}
    
    \item What MAP should be maintained to prevent secondary injury in SCI?
    
    MAP $>$ \underline{\hspace{3cm}}
    
    \vspace{0.3cm}
    
    \item How often should SCI patients be repositioned to prevent pressure injuries?
    
    Every \underline{\hspace{3cm}} hours
    
    \vspace{0.3cm}
    
    \item What medication is given for bradycardia in neurogenic shock?
    
    \underline{\hspace{6cm}}
    
\end{enumerate}

\newpage
\section*{ANSWER KEY - Sections 12-19}

\subsection*{Section 12: SBAR Communication}
\textbf{S:} T4; headache; 178/102; 100/70

\textbf{B:} 3; motorcycle; 2; bladder training; nipple line

\textbf{A:} autonomic dysreflexia; bladder distension

\textbf{R:} upright; Foley; nifedipine

\subsection*{Section 13: Clinical Judgment Scenarios}

\textbf{Scenario A:}
\begin{enumerate}
    \item Tension pneumothorax
    \item a) Tracheal deviation AWAY from affected side (to the right),
b) Absent breath sounds on the left
    
    c) Hypotension (BP 88/56)
    
    d) JVD (distended neck veins)
    
    e) Tachycardia (HR 128)
    
    \item Needle decompression at the 2nd intercostal space, midclavicular line on the LEFT side
    
    \item Airway with C-spine protection (A) - because the patient has neck pain and potential spinal injury from high-speed collision
\end{enumerate}

\textbf{Scenario B:}
\begin{enumerate}
    \item Class III Hemorrhagic Shock
    
    \item Evidence:
    \begin{itemize}
        \item HR 132 (Class III = 120-140)
        \item BP 86/54 - Decreased (Class III)
        \item Mental status: Confused (Class III)
        \item Estimated blood loss from pelvic fracture typically 1500-2000 mL
    \end{itemize}
    
    \item Interventions:
    \begin{itemize}
        \item Two large-bore IVs (16-gauge or larger)
        \item Warm crystalloid fluid resuscitation (LR preferred)
        \item Activate Massive Transfusion Protocol (MTP)
        \item Permissive hypotension (target SBP 80-90) until surgical control
        \item Pelvic binder application
        \item Prepare for OR/interventional radiology
    \end{itemize}
\end{enumerate}

\textbf{Scenario C:}
\begin{enumerate}
    \item Autonomic Dysreflexia
    
    \item Rationale: T6 injury is AT the T6 level, and autonomic dysreflexia occurs in SCI at or above T6. The classic presentation includes:
    \begin{itemize}
        \item Sudden severe hypertension (BP 192/108, baseline was likely much lower)
        \item Pounding headache
        \item Flushing and sweating ABOVE the level of injury
        \item Bradycardia (HR 52)
        \item Triggered by noxious stimulus below injury level (full bladder from clamped catheter)
    \end{itemize}
    
    \item Priority interventions in order:
    \begin{enumerate}[label=\alph*)]
        \item Sit the patient upright immediately (uses orthostatic effect to lower BP)
        \item Unclamp the Foley catheter to empty the bladder
        \item Loosen any restrictive clothing
        \item Monitor BP every 5 minutes
        \item Administer nifedipine if BP remains elevated
        \item Notify provider immediately
    \end{enumerate}
\end{enumerate}

\subsection*{Section 14: Triage Color Matching}
\begin{enumerate}
    \item RED - Immediate
    \item YELLOW - Delayed
    \item GREEN - Minor
    \item BLACK - Expectant
    \item RED - Immediate
    \item BLACK - Expectant
    \item YELLOW - Delayed
    \item GREEN - Minor
\end{enumerate}

\subsection*{Section 15: SCI Level Matching}
\begin{enumerate}
    \item C2 $\rightarrow$ A (Ventilator-dependent for life)
    \item C6 $\rightarrow$ C (Quadriplegia with some arm function)
    \item T8 $\rightarrow$ B (Paraplegia with full arm use)
    \item L3 $\rightarrow$ D (Paraplegia with variable leg function)
    \item S2 $\rightarrow$ E (Bowel/bladder dysfunction)
\end{enumerate}

\subsection*{Section 16: Beck's Triad}
The three components of Beck's Triad for cardiac tamponade:
\begin{enumerate}
    \item JVD (Jugular Vein Distension)
    \item Hypotension
    \item Muffled heart sounds
\end{enumerate}

Treatment: Pericardiocentesis

\subsection*{Section 17: Trauma Triad of Death}
\begin{enumerate}
    \item Hypothermia
    \item Acidosis
    \item Coagulopathy
\end{enumerate}

Prevention strategies:
\begin{itemize}
    \item Use warm IV fluids and blood products
    \item Apply warm blankets, use warming devices
    \item Limit exposure time during assessment
    \item Correct acidosis with adequate resuscitation
    \item Early activation of MTP to prevent dilutional coagulopathy
\end{itemize}

\subsection*{Section 18: True/False Answers}
\begin{enumerate}
    \item TRUE - Always assume C-spine injury in all trauma patients until ruled out
    \item FALSE - Tension pneumothorax causes tracheal deviation AWAY from the affected side (the pressure pushes structures away)
    \item TRUE - Neurogenic shock presents with hypotension AND bradycardia due to loss of sympathetic tone (cannot mount compensatory tachycardia)
    \item FALSE - Autonomic dysreflexia occurs in SCI at or ABOVE T6 (not below)
    \item TRUE - Permissive hypotension (SBP 80-90) is used in penetrating trauma to prevent ``popping the clot'' before surgical control
    \item FALSE - In START triage, patients who can walk are tagged GREEN (Minor), not YELLOW
    \item TRUE - C3, C4, C5 keeps the diaphragm alive - injury above C3 requires mechanical ventilation
    \item FALSE - Spinal shock has NORMAL hemodynamics; NEUROGENIC shock causes hypotension with bradycardia
    \item TRUE - The first step in autonomic dysreflexia is to sit the patient upright to use gravity/orthostatic effect to lower BP
    \item TRUE - MTP ratio is 1:1:1 (packed RBCs : FFP : Platelets) to prevent dilutional coagulopathy
\end{enumerate}

\subsection*{Section 19: Quick Recall Answers}
\begin{enumerate}
    \item ABCDE:
    \begin{itemize}
        \item A = Airway with cervical spine protection
        \item B = Breathing and ventilation
        \item C = Circulation with hemorrhage control
        \item D = Disability (neurological status - GCS, pupils)
        \item E = Exposure/Environmental control
    \end{itemize}
    
    \item AMPLE:
    \begin{itemize}
        \item A = Allergies
        \item M = Medications
        \item P = Past medical history
        \item L = Last meal
        \item E = Events leading to injury
    \end{itemize}
    
    \item START = Simple Triage and Rapid Treatment
    
    \item MTP ratio: 1 : 1 : 1
    
    \item IV catheter: 16-gauge or larger, quantity: 2 (two large-bore IVs)
    
    \item Preferred crystalloid: Lactated Ringer's (LR) - warmed
    
    \item Massive hemothorax: $>$ 1500 mL
    
    \item MAP for SCI: MAP $>$ 85
    
    \item Repositioning frequency: Every 2 hours
    
    \item Medication for bradycardia in neurogenic shock: Atropine
\end{enumerate}

\newpage
\section*{Section 20: Comprehensive Case Study Practice}

\begin{tcolorbox}[colback=red!5!white,colframe=red!75!black,title=NGN Unfolding Case Study - Multi-System Trauma]

\textbf{Initial Presentation:}

A 28-year-old male is brought to the ED by EMS after a high-speed motorcycle collision at approximately 60 mph. He was wearing a helmet. EMS reports he was found unconscious at the scene; he regained consciousness but is confused.

\textbf{On arrival assessment:}
\begin{itemize}
    \item GCS 12 (Eyes 3, Verbal 3, Motor 6)
    \item BP 88/56 mmHg
    \item HR 128 bpm
    \item RR 30
    \item SpO2 90\% on 15L non-rebreather mask
    \item Neck pain and bilateral leg numbness
    \item Trachea deviated to the RIGHT
    \item Absent breath sounds on the LEFT
    \item Distended neck veins (JVD)
    \item Rigid abdomen
    \item Pelvis unstable on compression
    \item 2.5 cm laceration on forehead
\end{itemize}

\end{tcolorbox}

\subsection*{Case Study Questions}

\textbf{Question 1: Recognition of Cues}

List ALL abnormal findings from this assessment and categorize them by body system:

\vspace{0.3cm}
\textbf{Neurological:}

\underline{\hspace{12cm}}

\underline{\hspace{12cm}}

\vspace{0.3cm}
\textbf{Respiratory:}

\underline{\hspace{12cm}}

\underline{\hspace{12cm}}

\vspace{0.3cm}
\textbf{Cardiovascular:}

\underline{\hspace{12cm}}

\underline{\hspace{12cm}}

\vspace{0.3cm}
\textbf{Abdominal/Pelvis:}

\underline{\hspace{12cm}}

\vspace{0.3cm}
\textbf{Integumentary:}

\underline{\hspace{12cm}}

\vspace{0.5cm}
\textbf{Question 2: Analysis of Cues}

Based on the respiratory findings (tracheal deviation to the right, absent breath sounds on the left, JVD, hypotension, tachypnea), what is the MOST likely diagnosis?

\underline{\hspace{10cm}}

What pathophysiology explains each of these findings?

\begin{itemize}
    \item Tracheal deviation to the RIGHT: \underline{\hspace{6cm}}
    \item Absent breath sounds LEFT: \underline{\hspace{6cm}}
    \item JVD: \underline{\hspace{6cm}}
    \item Hypotension: \underline{\hspace{6cm}}
\end{itemize}

\vspace{0.5cm}
\textbf{Question 3: Prioritization}

The patient has multiple injuries. Rank the following problems in order of priority (1 = highest priority):

\begin{tabular}{|c|p{10cm}|}
\hline
\textbf{Rank} & \textbf{Problem} \\
\hline
\underline{\hspace{1cm}} & Forehead laceration with active bleeding \\
\hline
\underline{\hspace{1cm}} & Tension pneumothorax \\
\hline
\underline{\hspace{1cm}} & Potential spinal cord injury \\
\hline
\underline{\hspace{1cm}} & Hemorrhagic shock from pelvic fracture \\
\hline
\underline{\hspace{1cm}} & Altered level of consciousness \\
\hline
\end{tabular}

\vspace{0.3cm}
Explain your rationale for the TOP priority:

\underline{\hspace{12cm}}

\underline{\hspace{12cm}}

\vspace{0.5cm}
\textbf{Question 4: Generate Solutions}

For the TOP priority problem identified above, what is the immediate intervention?

\underline{\hspace{12cm}}

Describe the procedure including:

\begin{itemize}
    \item Location: \underline{\hspace{8cm}}
    \item Equipment: \underline{\hspace{8cm}}
    \item Expected outcome: \underline{\hspace{8cm}}
\end{itemize}

\vspace{0.5cm}
\textbf{Question 5: Take Action}

After addressing the immediate life threat, list your next 5 interventions in order:

\begin{enumerate}
    \item \underline{\hspace{11cm}}
    \item \underline{\hspace{11cm}}
    \item \underline{\hspace{11cm}}
    \item \underline{\hspace{11cm}}
    \item \underline{\hspace{11cm}}
\end{enumerate}

\vspace{0.5cm}
\textbf{Question 6: Evaluate Outcomes}

After needle decompression and chest tube placement, what findings would indicate successful treatment of the tension pneumothorax?

\begin{itemize}
    \item Breath sounds: \underline{\hspace{8cm}}
    \item Trachea position: \underline{\hspace{8cm}}
    \item SpO2: \underline{\hspace{8cm}}
    \item Blood pressure: \underline{\hspace{8cm}}
    \item JVD: \underline{\hspace{8cm}}
\end{itemize}

\newpage
\subsection*{Case Study Answer Key}

\textbf{Question 1: Recognition of Cues}

\textbf{Neurological:}
\begin{itemize}
    \item GCS 12 (decreased from normal 15) - confusion
    \item Loss of consciousness at scene
    \item Neck pain
    \item Bilateral leg numbness (suggests spinal cord involvement)
\end{itemize}

\textbf{Respiratory:}
\begin{itemize}
    \item RR 30 (tachypnea - normal 12-20)
    \item SpO2 90\% on 15L (hypoxemia despite high-flow O2)
    \item Tracheal deviation to the RIGHT
    \item Absent breath sounds on the LEFT
    \item JVD (also cardiovascular)
\end{itemize}

\textbf{Cardiovascular:}
\begin{itemize}
    \item BP 88/56 (hypotension)
    \item HR 128 (tachycardia)
    \item JVD (distended neck veins)
\end{itemize}

\textbf{Abdominal/Pelvis:}
\begin{itemize}
    \item Rigid abdomen (suggests intra-abdominal bleeding)
    \item Pelvis unstable on compression (pelvic fracture - major blood loss)
\end{itemize}

\textbf{Integumentary:}
\begin{itemize}
    \item 2.5 cm forehead laceration
\end{itemize}

\vspace{0.5cm}
\textbf{Question 2: Analysis of Cues}

Most likely diagnosis: \textbf{LEFT-sided Tension Pneumothorax}

Pathophysiology:
\begin{itemize}
    \item Tracheal deviation to the RIGHT: Air trapped in left pleural space creates pressure that pushes mediastinal structures (including trachea) AWAY from affected side
    \item Absent breath sounds LEFT: Collapsed lung on left side cannot ventilate
    \item JVD: Increased intrathoracic pressure compresses vena cava, preventing venous return, causing backup into jugular veins
    \item Hypotension: Compressed vena cava decreases preload, reducing cardiac output
\end{itemize}

\vspace{0.5cm}
\textbf{Question 3: Prioritization}

\begin{enumerate}
    \item \textbf{Tension pneumothorax} (Breathing - immediate life threat)
    \item Potential spinal cord injury (Airway/C-spine - addressed simultaneously)
    \item Hemorrhagic shock from pelvic fracture (Circulation)
    \item Altered level of consciousness (Disability - assess after ABC)
    \item Forehead laceration (address after life threats)
\end{enumerate}

\textbf{Rationale:} Tension pneumothorax is immediately life-threatening because it causes cardiovascular collapse. The increasing pressure in the chest compresses the heart and great vessels, rapidly leading to death if not treated. This is a clinical diagnosis that requires immediate needle decompression - do NOT wait for imaging. Following ABCDE, after securing the airway with C-spine precautions, breathing emergencies take priority.

\vspace{0.5cm}
\textbf{Question 4: Generate Solutions}

Immediate intervention: \textbf{Needle decompression (needle thoracostomy)}

\begin{itemize}
    \item Location: 2nd intercostal space (ICS), midclavicular line (MCL), LEFT side - insert needle over the TOP of the 3rd rib to avoid the neurovascular bundle
    \item Equipment: Large-bore needle (14-16 gauge), at least 5 cm long
    \item Expected outcome: Rush of air heard as tension is released; improvement in vital signs, breath sounds, and oxygenation; followed by chest tube placement
\end{itemize}

\vspace{0.5cm}
\textbf{Question 5: Take Action}

After needle decompression:
\begin{enumerate}
    \item Maintain manual in-line cervical spine stabilization; apply rigid cervical collar
    \item Establish two large-bore IV access (16-gauge or larger)
    \item Initiate warm crystalloid fluid resuscitation (LR)
    \item Apply pelvic binder for unstable pelvis
    \item Activate Massive Transfusion Protocol (MTP)
    \item Prepare for chest tube insertion
    \item Assess GCS, pupils, motor function
    \item Apply direct pressure to forehead laceration
    \item Prepare for emergent OR/CT imaging once stabilized
\end{enumerate}

\vspace{0.5cm}
\textbf{Question 6: Evaluate Outcomes}

Successful treatment indicators:
\begin{itemize}
    \item Breath sounds: Present and equal bilaterally
    \item Trachea position: Midline
    \item SpO2: Improved ($>$94\%)
    \item Blood pressure: Improved (though may still be low from other injuries)
    \item JVD: Resolved (neck veins flat or normal)
    \item Respiratory rate: Decreased toward normal
    \item Patient comfort: Decreased respiratory distress
\end{itemize}

\newpage
\section*{Section 21: High-Yield Facts for Quick Review}

\begin{tcolorbox}[colback=yellow!10!white,colframe=yellow!75!black,title=MUST KNOW for Exam]

\textbf{Primary Survey Order:} A-B-C-D-E (Airway, Breathing, Circulation, Disability, Exposure)

\textbf{Tension Pneumothorax Signs:}
\begin{itemize}
    \item Tracheal deviation AWAY from affected side
    \item Absent breath sounds on affected side
    \item JVD
    \item Hypotension
    \item Treatment: Needle decompression at 2nd ICS, MCL - do NOT wait for X-ray!
\end{itemize}

\textbf{Beck's Triad (Cardiac Tamponade):}
\begin{itemize}
    \item JVD
    \item Hypotension
    \item Muffled heart sounds
    \item Treatment: Pericardiocentesis
\end{itemize}

\textbf{Trauma Triad of Death:}
\begin{itemize}
    \item Hypothermia
    \item Acidosis
    \item Coagulopathy
\end{itemize}

\textbf{Hemorrhagic Shock - Key Class Markers:}
\begin{itemize}
    \item Class I: Normal VS, normal mental status
    \item Class II: T
ach
ach
ycardia, anxiety, narrowed pulse pressure
    \item Class III: Hypotension, confusion, tachycardia 120-140
    \item Class IV: Severe hypotension, lethargy/unconscious, HR $>$140
\end{itemize}

\textbf{MTP Ratio:} 1:1:1 (PRBCs : FFP : Platelets)

\textbf{Neurogenic vs Hypovolemic Shock:}
\begin{itemize}
    \item Neurogenic: Bradycardia + Hypotension + Warm/dry skin
    \item Hypovolemic: Tachycardia + Hypotension + Cool/clammy skin
\end{itemize}

\textbf{Autonomic Dysreflexia:}
\begin{itemize}
    \item Occurs with SCI at T6 or above
    \item Triggered by noxious stimulus below injury (full bladder most common)
    \item Signs: Severe HTN, pounding headache, flushing/sweating above injury, bradycardia
    \item Treatment: Sit patient UP, remove stimulus, give nifedipine if needed
    \item EMERGENCY - can cause stroke!
\end{itemize}

\textbf{Spinal Shock:}
\begin{itemize}
    \item Temporary loss of all reflexes below injury
    \item Flaccid paralysis, areflexia
    \item Resolves when bulbocavernosus reflex returns
\end{itemize}

\textbf{SCI MAP Goal:} Maintain MAP $>$85 mmHg for 7 days

\textbf{Cushing's Triad (Increased ICP):}
\begin{itemize}
    \item Hypertension (widening pulse pressure)
    \item Bradycardia
    \item Irregular respirations
\end{itemize}

\textbf{GCS Scoring:}
\begin{itemize}
    \item Eye: 4 (spontaneous) to 1 (none)
    \item Verbal: 5 (oriented) to 1 (none)
    \item Motor: 6 (obeys commands) to 1 (none)
    \item Severe TBI: GCS $\leq$8 (intubate!)
\end{itemize}

\textbf{Burn Fluid Resuscitation (Parkland Formula):}
\begin{itemize}
    \item 4 mL $\times$ kg $\times$ \%TBSA
    \item Half in first 8 hours, half in next 16 hours
    \item Use LR (Lactated Ringer's)
\end{itemize}

\end{tcolorbox}

\newpage
\section*{Section 22: Comprehensive Practice Exam}

\begin{tcolorbox}[colback=gray!10!white,colframe=gray!75!black,title=Practice Exam Instructions]
This comprehensive exam covers all trauma and emergency nursing concepts. Select the BEST answer for each question. Time yourself - aim for 1-1.5 minutes per question.
\end{tcolorbox}

\vspace{0.5cm}

\textbf{1.} A patient arrives after a motor vehicle crash. The nurse notes tracheal deviation to the right, absent breath sounds on the left, and JVD. What is the priority intervention?

\begin{enumerate}[label=\Alph*.]
    \item Obtain a chest X-ray
    \item Perform needle decompression on the left side
    \item Establish IV access
    \item Administer oxygen via non-rebreather mask
\end{enumerate}

\vspace{0.3cm}

\textbf{2.} A patient with a T4 spinal cord injury suddenly develops a blood pressure of 210/120 mmHg, severe headache, and facial flushing. The nurse's FIRST action should be to:

\begin{enumerate}[label=\Alph*.]
    \item Administer prescribed antihypertensive medication
    \item Check the Foley catheter for obstruction
    \item Place the patient in Trendelenburg position
    \item Call a rapid response
\end{enumerate}

\vspace{0.3cm}

\textbf{3.} Which assessment finding differentiates neurogenic shock from hypovolemic shock?

\begin{enumerate}[label=\Alph*.]
    \item Hypotension
    \item Decreased urine output
    \item Bradycardia
    \item Altered mental status
\end{enumerate}

\vspace{0.3cm}

\textbf{4.} A trauma patient has an estimated blood loss of 1,800 mL. The nurse anticipates which class of hemorrhagic shock?

\begin{enumerate}[label=\Alph*.]
    \item Class I
    \item Class II
    \item Class III
    \item Class IV
\end{enumerate}

\vspace{0.3cm}

\textbf{5.} During the primary survey, a patient is found to have gurgling respirations. What is the nurse's priority action?

\begin{enumerate}[label=\Alph*.]
    \item Assess circulation
    \item Suction the airway
    \item Check pupil response
    \item Apply a cervical collar
\end{enumerate}

\vspace{0.3cm}

\textbf{6.} The nurse is caring for a patient with suspected cardiac tamponade. Which findings would the nurse expect? (Select all that apply)

\begin{enumerate}[label=\Alph*.]
    \item Muffled heart sounds
    \item Tracheal deviation
    \item JVD
    \item Hypotension
    \item Hyperresonance to percussion
\end{enumerate}

\vspace{0.3cm}

\textbf{7.} A patient with a C5 spinal cord injury is at risk for respiratory compromise because:

\begin{enumerate}[label=\Alph*.]
    \item The phrenic nerve innervating the diaphragm originates at C3-C5
    \item Cervical injuries always require intubation
    \item The patient will develop autonomic dysreflexia
    \item Spinal shock causes permanent respiratory paralysis
\end{enumerate}

\vspace{0.3cm}

\textbf{8.} When using the START triage system, a patient who is not breathing after repositioning the airway should be tagged:

\begin{enumerate}[label=\Alph*.]
    \item Red (Immediate)
    \item Yellow (Delayed)
    \item Green (Minor)
    \item Black (Expectant)
\end{enumerate}

\vspace{0.3cm}

\textbf{9.} A patient with 40\% TBSA burns weighing 70 kg requires fluid resuscitation. Using the Parkland formula, how much fluid should be administered in the first 8 hours?

\begin{enumerate}[label=\Alph*.]
    \item 5,600 mL
    \item 11,200 mL
    \item 2,800 mL
    \item 8,400 mL
\end{enumerate}

\vspace{0.3cm}

\textbf{10.} The nurse is preparing to receive a trauma patient. Which IV access is most appropriate?

\begin{enumerate}[label=\Alph*.]
    \item One 22-gauge IV in the hand
    \item Two 18-gauge IVs in the antecubital fossae
    \item One 20-gauge IV with a saline lock
    \item Two 16-gauge or larger IVs
\end{enumerate}

\vspace{0.3cm}

\textbf{11.} A patient presents with Cushing's triad. The nurse recognizes this indicates:

\begin{enumerate}[label=\Alph*.]
    \item Cardiac tamponade
    \item Increased intracranial pressure
    \item Tension pneumothorax
    \item Neurogenic shock
\end{enumerate}

\vspace{0.3cm}

\textbf{12.} Which GCS score indicates severe traumatic brain injury requiring immediate airway management?

\begin{enumerate}[label=\Alph*.]
    \item GCS 15
    \item GCS 12
    \item GCS 9
    \item GCS 7
\end{enumerate}

\vspace{0.3cm}

\textbf{13.} The nurse is caring for a patient in spinal shock. Which finding indicates resolution of spinal shock?

\begin{enumerate}[label=\Alph*.]
    \item Return of blood pressure to normal
    \item Presence of bulbocavernosus reflex
    \item Ability to move extremities
    \item Normal body temperature
\end{enumerate}

\vspace{0.3cm}

\textbf{14.} During massive transfusion protocol, the ratio of PRBCs to FFP to platelets should be:

\begin{enumerate}[label=\Alph*.]
    \item 2:1:1
    \item 1:1:1
    \item 3:1:1
    \item 1:2:1
\end{enumerate}

\vspace{0.3cm}

\textbf{15.} A patient with a flail chest is at risk for respiratory failure because:

\begin{enumerate}[label=\Alph*.]
    \item The trachea deviates toward the affected side
    \item Paradoxical chest wall movement impairs ventilation
    \item Air accumulates in the pleural space
    \item Blood fills the pericardial sac
\end{enumerate}

\vspace{0.3cm}

\textbf{16.} The preferred crystalloid fluid for trauma resuscitation is:

\begin{enumerate}[label=\Alph*.]
    \item Normal saline (0.9\% NaCl)
    \item D5W
    \item Lactated Ringer's
    \item D5 1/2 NS
\end{enumerate}

\vspace{0.3cm}

\textbf{17.} A patient with a SCI at T6 is at risk for autonomic dysreflexia. Which nursing intervention is essential?

\begin{enumerate}[label=\Alph*.]
    \item Maintain the patient in supine position
    \item Ensure regular bladder emptying
    \item Restrict fluid intake
    \item Apply compression stockings continuously
\end{enumerate}

\vspace{0.3cm}

\textbf{18.} Massive hemothorax is defined as blood loss greater than:

\begin{enumerate}[label=\Alph*.]
    \item 500 mL
    \item 1,000 mL
    \item 1,500 mL
    \item 2,000 mL
\end{enumerate}

\vspace{0.3cm}

\textbf{19.} The nurse is assessing a patient using the AMPLE history. What does the "P" stand for?

\begin{enumerate}[label=\Alph*.]
    \item Pain level
    \item Past medical history/Pregnancy
    \item Pulse rate
    \item Primary complaint
\end{enumerate}

\vspace{0.3cm}

\textbf{20.} A patient with suspected spinal cord injury should have MAP maintained at:

\begin{enumerate}[label=\Alph*.]
    \item $>$60 mmHg
    \item $>$70 mmHg
    \item $>$85 mmHg
    \item $>$100 mmHg
\end{enumerate}

\newpage
\section*{Practice Exam Answer Key}

\begin{enumerate}
    \item \textbf{B} - Needle decompression on the left side. Tension pneumothorax is a life-threatening emergency that requires immediate intervention. Do NOT wait for X-ray confirmation.
    
    \item \textbf{B} - Check the Foley catheter for obstruction. The first action in autonomic dysreflexia is to identify and remove the triggering stimulus. Bladder distension is the most common cause.
    
    \item \textbf{C} - Bradycardia. Neurogenic shock presents with bradycardia due to loss of sympathetic tone, while hypovolemic shock presents with tachycardia as a compensatory mechanism.
    
    \item \textbf{C} - Class III. Class III hemorrhagic shock involves 1,500-2,000 mL (30-40\%) blood loss with hypotension and confusion.
    
    \item \textbf{B} - Suction the airway. Airway is the first priority in the primary survey. Gurgling indicates fluid/secretions that must be cleared.
    
    \item \textbf{A, C, D} - Beck's triad includes muffled heart sounds, JVD, and hypotension. Tracheal deviation and hyperresonance are associated with tension pneumothorax.
    
    \item \textbf{A} - The phrenic nerve originates at C3-C5. Injuries at or above this level can impair diaphragmatic function and breathing.
    
    \item \textbf{D} - Black (Expectant). In START triage, patients who are not breathing even after airway repositioning are tagged black.
    
    \item \textbf{A} - 5,600 mL. Parkland: 4 $\times$ 70 $\times$ 40 = 11,200 mL total. Half in first 8 hours = 5,600 mL.
    
    \item \textbf{D} - Two 16-gauge or larger IVs. Trauma patients require large-bore access for rapid fluid and blood product administration.
    
    \item \textbf{B} - Increased intracranial pressure. Cushing's triad (hypertension, bradycardia, irregular respirations) indicates dangerously elevated ICP.
    
    \item \textbf{D} - GCS 7. Severe TBI is defined as GCS $\leq$8, requiring immediate airway protection/intubation.
    
    \item \textbf{B} - Presence of bulbocavernosus reflex. Return of this reflex indicates the end of spinal shock.
    
    \item \textbf{B} - 1:1:1. This balanced ratio helps prevent the trauma triad of death.
    
    \item \textbf{B} - Paradoxical chest wall movement impairs ventilation. The flail segment moves opposite to the rest of the chest wall during breathing.
    
    \item \textbf{C} - Lactated Ringer's. LR is preferred because it is more physiologically similar to plasma and helps buffer acidosis.
    
    \item \textbf{B} - Ensure regular bladder emptying. Bladder distension is the most common trigger for autonomic dysreflexia.
    
    \item \textbf{C} - 1,500 mL. Massive hemothorax is defined as $>$1,500 mL of blood in the pleural space.
    
    \item \textbf{B} - Past medical history/Pregnancy. AMPLE = Allergies, Medications, Past medical history/Pregnancy, Last meal, Events leading to injury.
    
    \item \textbf{C} - $>$85 mmHg. Current guidelines recommend maintaining MAP $>$85 mmHg for 7 days to prevent secondary spinal cord injury.
\end{enumerate}

\newpage
\section*{Section 23: Critical Thinking Case Studies}

\begin{tcolorbox}[colback=purple!5!white,colframe=purple!75!black,title=Complex Case Study 1: Multi-System Trauma]

\textbf{Scenario:} A 28-year-old male is brought to the ED after being ejected from a vehicle during a rollover crash. He was found 20 feet from the vehicle, unresponsive at the scene. EMS reports GCS of 7 (E2V1M4), BP 82/50, HR 134, RR 28, SpO2 88\% on 15L non-rebreather.

\textbf{On arrival:}
\begin{itemize}
    \item Airway: Snoring respirations, blood in oropharynx
    \item Breathing: Decreased breath sounds on right, subcutaneous emphysema right chest
    \item Circulation: Cool, pale, diaphoretic; weak radial pulses; obvious open femur fracture right leg with active bleeding
    \item Disability: GCS 7, pupils 4mm equal and reactive
    \item Exposure: Multiple abrasions, rigid abdomen
\end{itemize}
\end{tcolorbox}

\textbf{Questions:}

\textbf{1.} List the immediate life threats identified in this patient in order of priority:

\vspace{2cm}

\textbf{2.} What is the FIRST intervention the trauma team should perform?

\vspace{1.5cm}

\textbf{3.} Based on the chest findings, what injury should be suspected and what intervention is needed?

\vspace{1.5cm}

\textbf{4.} Calculate the estimated blood loss class based on the vital signs and clinical presentation:

\vspace{1.5cm}

\textbf{5.} What does the rigid abdomen suggest, and what diagnostic test would confirm this?

\vspace{1.5cm}

\textbf{6.} Develop a prioritized nursing care plan for the first 30 minutes:

\vspace{3cm}

\begin{tcolorbox}[colback=green!5!white,colframe=green!75!black,title=Case Study 1 - Answers]

\textbf{1. Life threats in priority order:}
\begin{enumerate}
    \item Compromised airway (snoring, blood in oropharynx, GCS 7)
    \item Right tension pneumothorax or hemopneumothorax (decreased breath sounds, subcutaneous emphysema)
    \item Hemorrhagic shock (hypotension, tachycardia, cool/pale skin, open femur fracture)
    \item Potential intra-abdominal hemorrhage (rigid abdomen)
    \item Traumatic brain injury (GCS 7)
\end{enumerate}

\textbf{2. FIRST intervention:}
Secure the airway with rapid sequence intubation (RSI) while maintaining cervical spine immobilization. GCS $\leq$8 requires definitive airway management.

\textbf{3. Chest injury and intervention:}
Suspected: Tension pneumothorax or hemopneumothorax on the right side.
Intervention: Needle decompression at 2nd ICS, MCL on the RIGHT side, followed by chest tube insertion.

\textbf{4. Hemorrhagic shock class:}
Class III-IV hemorrhagic shock
\begin{itemize}
    \item HR 134 (Class III-IV range)
    \item BP 82/50 (significantly hypotensive)
    \item RR 28 (elevated)
    \item Cool, pale, diaphoretic (Class III-IV)
    \item Open femur fracture can cause 1,500-2,000 mL blood loss alone
\end{itemize}

\textbf{5. Rigid abdomen:}
Suggests intra-abdominal hemorrhage, likely from solid organ injury (spleen, liver) or hollow viscus perforation.
Diagnostic test: FAST exam (Focused Assessment with Sonography for Trauma) at bedside; CT abdomen/pelvis if patient stable enough.

\textbf{6. Prioritized nursing care plan (first 30 minutes):}
\begin{enumerate}
    \item Minutes 0-5: Assist with RSI, maintain C-spine immobilization, confirm ETT placement
    \item Minutes 5-10: Assist with needle decompression/chest tube; establish two 16G IV access
    \item Minutes 10-15: Initiate MTP; hang O-negative blood; apply direct pressure to femur wound; apply tourniquet if needed
    \item Minutes 15-20: FAST exam; apply pelvic binder; warm
fluids; prevent hypothermia with warm blankets/Bair Hugger
    \item Minutes 20-25: Continuous VS monitoring; Foley catheter (if no contraindication); NG tube
    \item Minutes 25-30: Prepare for emergent OR or CT; document all interventions; update trauma team
\end{enumerate}
\end{tcolorbox}

\newpage
\subsection*{Case Study 2: Spinal Cord Injury with Neurogenic Shock}

\textbf{Scenario:}

A 22-year-old male was diving into a shallow pool and struck his head on the bottom. Bystanders pulled him from the water and called 911. Upon EMS arrival, the patient was conscious but unable to move his arms or legs. He was transported with full spinal immobilization.

\textbf{ED Assessment Findings:}
\begin{itemize}
    \item Alert and oriented, GCS 15
    \item C/O neck pain, unable to feel or move extremities
    \item BP 78/52, HR 54, RR 18 (shallow), SpO2 94\% on room air
    \item Skin warm and dry, flushed below clavicles
    \item Flaccid paralysis of all extremities
    \item No response to painful stimuli below clavicles
    \item Priapism noted
    \item Abdomen soft, non-tender, bowel sounds absent
\end{itemize}

\vspace{0.5cm}
\textbf{Answer the following questions:}

\textbf{1.} What type of shock is this patient experiencing? How do you differentiate it from hypovolemic shock?

\vspace{1.5cm}

\textbf{2.} At what spinal cord level do you suspect the injury? What findings support this?

\vspace{1.5cm}

\textbf{3.} Why is this patient at high risk for respiratory failure?

\vspace{1.5cm}

\textbf{4.} What are the priority nursing interventions?

\vspace{1.5cm}

\textbf{5.} What complication should you monitor for once the patient is admitted to the ICU, especially if the injury is at T6 or above?

\vspace{1.5cm}

\textbf{6.} The patient's family asks about prognosis. What information can you provide about spinal shock vs. permanent injury?

\vspace{2cm}

\begin{tcolorbox}[colback=green!5!white,colframe=green!75!black,title=Case Study 2 - Answers]

\textbf{1. Type of shock and differentiation:}

\textbf{Neurogenic shock} - caused by loss of sympathetic tone due to spinal cord injury.

\begin{center}
\begin{tabular}{|l|l|l|}
\hline
\textbf{Finding} & \textbf{Neurogenic Shock} & \textbf{Hypovolemic Shock} \\
\hline
Heart rate & Bradycardia & Tachycardia \\
Blood pressure & Hypotension & Hypotension \\
Skin & Warm, dry, flushed & Cool, clammy, pale \\
Mechanism & Loss of sympathetic tone & Blood/fluid loss \\
\hline
\end{tabular}
\end{center}

\textbf{2. Suspected spinal cord level:}

\textbf{High cervical injury (C3-C5)} based on:
\begin{itemize}
    \item Complete quadriplegia (no movement of arms or legs)
    \item Loss of sensation below clavicles
    \item Shallow respirations (diaphragm innervated by C3-C5)
    \item Priapism (loss of sympathetic control)
    \item Neurogenic shock (injury above T6)
\end{itemize}

\textbf{3. Risk for respiratory failure:}

\begin{itemize}
    \item Diaphragm innervation (C3-C5) may be partially or fully compromised
    \item Loss of intercostal muscle function (T1-T12)
    \item Loss of abdominal muscle function (cannot cough effectively)
    \item Shallow breathing pattern observed
    \item SpO2 already decreased (94\%)
    \item At risk for ascending edema worsening respiratory function
\end{itemize}

\textbf{4. Priority nursing interventions:}

\begin{enumerate}
    \item Maintain strict spinal immobilization
    \item Supplemental oxygen; prepare for possible intubation
    \item IV access with warm crystalloid fluids (cautious - avoid fluid overload)
    \item Vasopressors (norepinephrine or phenylephrine) for neurogenic shock
    \item Atropine available for symptomatic bradycardia
    \item Foley catheter for urinary retention
    \item NG tube for gastric decompression (paralytic ileus)
    \item Temperature management (poikilothermia - cannot regulate temperature)
    \item Continuous cardiac and respiratory monitoring
    \item Prepare for emergent MRI
\end{enumerate}

\textbf{5. Complication to monitor (injury T6 or above):}

\textbf{Autonomic Dysreflexia} - though this typically occurs after spinal shock resolves (days to weeks post-injury). Monitor for:
\begin{itemize}
    \item Severe hypertension (SBP can exceed 200 mmHg)
    \item Pounding headache
    \item Flushing and sweating above level of injury
    \item Bradycardia
    \item Nasal congestion
\end{itemize}
Common triggers: Full bladder (most common), full bowel, skin irritation, tight clothing.

\textbf{6. Information for family regarding prognosis:}

\begin{itemize}
    \item \textbf{Spinal shock} is a temporary condition lasting hours to weeks where all reflexes below the injury are absent - this does NOT predict permanent outcome
    \item Cannot determine permanent neurological status until spinal shock resolves
    \item Resolution of spinal shock is indicated by return of reflexes (bulbocavernosus reflex)
    \item After spinal shock resolves, remaining deficits are more likely permanent
    \item Early surgical decompression (if indicated) and rehabilitation can optimize outcomes
    \item Provide emotional support; involve social work, chaplain services
\end{itemize}
\end{tcolorbox}

\newpage
\subsection*{Case Study 3: Penetrating Abdominal Trauma}

\textbf{Scenario:}

A 28-year-old male arrives via EMS after being stabbed in the left upper quadrant of the abdomen during an altercation. He is alert but anxious and complaining of severe abdominal pain.

\textbf{Assessment Findings:}
\begin{itemize}
    \item Alert, oriented, GCS 15, anxious
    \item BP 102/68, HR 118, RR 24, SpO2 97\% on room air
    \item Stab wound to LUQ, approximately 3 cm, minimal external bleeding
    \item Abdomen rigid, diffusely tender, guarding present
    \item Bowel sounds absent
    \item Skin cool and diaphoretic
    \item Pain rated 10/10
\end{itemize}

\vspace{0.5cm}
\textbf{Answer the following questions:}

\textbf{1.} What organ is most likely injured based on the wound location?

\vspace{1.5cm}

\textbf{2.} Why is minimal external bleeding NOT reassuring in this case?

\vspace{1.5cm}

\textbf{3.} What class of hemorrhagic shock is this patient in?

\vspace{1.5cm}

\textbf{4.} Should you remove or explore the wound? Why or why not?

\vspace{1.5cm}

\textbf{5.} What diagnostic studies would you anticipate?

\vspace{1.5cm}

\textbf{6.} Develop an SBAR handoff for this patient going to the OR.

\vspace{3cm}

\begin{tcolorbox}[colback=green!5!white,colframe=green!75!black,title=Case Study 3 - Answers]

\textbf{1. Most likely injured organ:}

\textbf{Spleen} - located in the left upper quadrant, highly vascular, and commonly injured in LUQ penetrating trauma. Other possibilities include:
\begin{itemize}
    \item Left kidney
    \item Stomach
    \item Splenic flexure of colon
    \item Left hemidiaphragm
    \item Tail of pancreas
\end{itemize}

\textbf{2. Why minimal external bleeding is NOT reassuring:}

\begin{itemize}
    \item External bleeding does not correlate with internal bleeding
    \item The spleen is highly vascular and can bleed massively into the peritoneal cavity
    \item Rigid abdomen and absent bowel sounds suggest significant intra-abdominal pathology
    \item Patient is already showing signs of shock (tachycardia, cool/diaphoretic skin)
    \item ``The bleeding you can't see is the bleeding that kills''
\end{itemize}

\textbf{3. Class of hemorrhagic shock:}

\textbf{Class II hemorrhagic shock} (estimated 750-1500 mL blood loss):
\begin{itemize}
    \item HR 118 (100-120 range)
    \item BP 102/68 - narrowed pulse pressure (102-68=34)
    \item RR 24 (elevated)
    \item Mental status: Anxious (Class II finding)
    \item Cool, diaphoretic skin
\end{itemize}
Note: Patient is at risk for rapid progression to Class III.

\textbf{4. Should you remove or explore the wound?}

\textbf{NO!} Never remove impaled objects or explore penetrating wounds in the ED.
\begin{itemize}
    \item Exploration should only occur in the OR with surgical control
    \item Removing objects or probing can dislodge clots and worsen bleeding
    \item Cover the wound with sterile saline-moistened gauze
    \item Stabilize any impaled objects in place
\end{itemize}

\textbf{5. Anticipated diagnostic studies:}

\begin{itemize}
    \item FAST exam - quick bedside assessment for free fluid
    \item Type and crossmatch - prepare for transfusion
    \item CBC, CMP, coagulation studies, lactate
    \item If stable: CT abdomen/pelvis with IV contrast
    \item If unstable: Direct to OR for exploratory laparotomy (no CT)
    \item Chest X-ray - rule out diaphragm injury, hemothorax
\end{itemize}

\textbf{6. SBAR Handoff to OR:}

\begin{tcolorbox}[colback=blue!5!white,colframe=blue!75!black,title=SBAR Handoff]
\textbf{S - Situation:}
``I'm calling about Mr. J, a 28-year-old male with a penetrating stab wound to the left upper quadrant who needs emergent exploratory laparotomy.''

\textbf{B - Background:}
``He was stabbed approximately 45 minutes ago during an altercation. He has a 3 cm stab wound to the LUQ. No significant medical history. Last ate 3 hours ago. Allergies: NKDA.''

\textbf{A - Assessment:}
``He is in Class II hemorrhagic shock with HR 118, BP 102/68, rigid abdomen, and absent bowel sounds. FAST exam positive for free fluid in the LUQ. I suspect splenic injury with ongoing intra-abdominal hemorrhage.''

\textbf{R - Recommendation:}
``He needs emergent surgical exploration. We have 2 large-bore IVs, type and cross for 6 units sent, MTP on standby. He's received 1 liter of LR. Ready for transport to OR when you're ready.''
\end{tcolorbox}

\end{tcolorbox}

\newpage
\section{Final Comprehensive Practice Examination}

\begin{center}
\textbf{\Large Trauma Nursing Comprehensive Final Exam}\\
\vspace{0.3cm}
\textit{75 Questions - NCLEX-Style and NGN Format}\\
\textit{Time Limit: 90 minutes}
\end{center}

\vspace{0.5cm}

\textbf{Instructions:} This examination covers all content from the trauma nursing study guide. Questions are formatted in NCLEX-style (multiple choice, select all that apply) and Next Generation NCLEX formats (case studies, bowtie, matrix). Choose the best answer(s) for each question.

\vspace{0.5cm}

\subsection*{Questions 1-15: Primary Survey and Initial Assessment}

\begin{enumerate}

\item A trauma patient arrives with gurgling respirations, blood in the oropharynx, and a GCS of 6. What is the nurse's FIRST action?

\begin{enumerate}[label=\Alph*.]
    \item Obtain a chest X-ray
    \item Establish IV access
    \item Suction the airway and prepare for intubation
    \item Assess pupil response
\end{enumerate}

\item The nurse is caring for a patient with suspected cervical spine injury. Which intervention is MOST important during airway management?

\begin{enumerate}[label=\Alph*.]
    \item Hyperextend the neck to open the airway
    \item Use the head-tilt chin-lift maneuver
    \item Maintain manual in-line stabilization during intubation
    \item Remove the cervical collar to visualize the airway
\end{enumerate}

\item A trauma patient has the following findings: absent breath sounds on the right, tracheal deviation to the left, JVD, and hypotension. What condition does the nurse suspect?

\begin{enumerate}[label=\Alph*.]
    \item Cardiac tamponade
    \item Right tension pneumothorax
    \item Massive hemothorax
    \item Flail chest
\end{enumerate}

\item For the patient in question 3, what is the immediate treatment?

\begin{enumerate}[label=\Alph*.]
    \item Pericardiocentesis
    \item Needle decompression at 2nd ICS, right midclavicular line
    \item Chest tube insertion at 5th ICS
    \item Emergency thoracotomy
\end{enumerate}

\item A patient involved in a high-speed MVC has a GCS of 7 (E2, V2, M3). The nurse knows this patient requires:

\begin{enumerate}[label=\Alph*.]
    \item Supplemental oxygen via nasal cannula
    \item Non-rebreather mask at 15 L/min
    \item Definitive airway management (intubation)
    \item Bag-valve-mask ventilation only
\end{enumerate}

\item \textbf{Select all that apply.} Which findings indicate adequate breathing in a trauma patient?

\begin{enumerate}[label=\Alph*.]
    \item Respiratory rate of 14 breaths/min
    \item SpO2 of 98\% on room air
    \item Use of accessory muscles
    \item Equal bilateral chest rise
    \item Symmetrical breath sounds
    \item Paradoxical chest movement
\end{enumerate}

\item The nurse is assessing circulation in a trauma patient. Which finding is MOST concerning for hemorrhagic shock?

\begin{enumerate}[label=\Alph*.]
    \item Capillary refill of 2 seconds
    \item Blood pressure of 118/76 mmHg
    \item Heart rate of 128 bpm with weak, thready pulse
    \item Warm, dry skin
\end{enumerate}

\item During the primary survey, the nurse notes that a patient withdraws from painful stimuli but does not open eyes or speak. What is this patient's GCS score?

\begin{enumerate}[label=\Alph*.]
    \item 4
    \item 5
    \item 6
    \item 7
\end{enumerate}

\item A trauma patient has unequal pupils with the right pupil fixed and dilated. This finding suggests:

\begin{enumerate}[label=\Alph*.]
    \item Bilateral cerebral contusion
    \item Uncal herniation with compression of CN III on the right
    \item Optic nerve injury
    \item Normal variation
\end{enumerate}

\item The ``E'' in the primary survey stands for Exposure/Environment. Which nursing action is MOST important during this phase?

\begin{enumerate}[label=\Alph*.]
    \item Obtain a detailed medical history
    \item Remove all clothing and logroll to assess posterior surfaces
    \item Apply warm blankets and prevent hypothermia
    \item Both B and C
\end{enumerate}

\item \textbf{Select all that apply.} Which patients require cervical spine immobilization?

\begin{enumerate}[label=\Alph*.]
    \item Alert patient with neck pain after MVC
    \item Intoxicated patient after fall from standing
    \item Patient with focal neurological deficits after diving injury
    \item Elderly patient with head laceration after mechanical fall
    \item Patient with distracting injury (open femur fracture) after MVC
\end{enumerate}

\item A patient has a blood pressure of 88/62 mmHg and heart rate of 52 bpm after a motorcycle crash with suspected spinal cord injury. The nurse recognizes this as:

\begin{enumerate}[label=\Alph*.]
Hypovolemic shock
    \item Neurogenic shock
    \item Cardiogenic shock
    \item Septic shock
\end{enumerate}

\item The nurse is caring for a patient with suspected tension pneumothorax. Which intervention takes priority?

\begin{enumerate}[label=\Alph*.]
    \item Obtain a chest X-ray
    \item Prepare for needle decompression
    \item Administer high-flow oxygen
    \item Establish IV access
\end{enumerate}

\item A patient with multiple trauma has the following vital signs: HR 142, BP 78/40, RR 32, confused and agitated. The nurse estimates this patient is in which class of hemorrhagic shock?

\begin{enumerate}[label=\Alph*.]
    \item Class I
    \item Class II
    \item Class III
    \item Class IV
\end{enumerate}

\item The Massive Transfusion Protocol (MTP) calls for blood products in what ratio?

\begin{enumerate}[label=\Alph*.]
    \item 3:1:1 (PRBCs : FFP : Platelets)
    \item 1:1:1 (PRBCs : FFP : Platelets)
    \item 2:1:1 (PRBCs : FFP : Platelets)
    \item 1:2:1 (PRBCs : FFP : Platelets)
\end{enumerate}

\item \textbf{Select all that apply.} Which findings indicate successful treatment of tension pneumothorax after needle decompression?

\begin{enumerate}[label=\Alph*.]
    \item Rush of air heard upon needle insertion
    \item Improved oxygen saturation
    \item Trachea returning to midline
    \item Decreased heart rate
    \item Resolution of JVD
    \item Bilateral breath sounds
\end{enumerate}

\item A trauma patient has Beck's triad. The nurse anticipates preparing for which procedure?

\begin{enumerate}[label=\Alph*.]
    \item Chest tube insertion
    \item Needle decompression
    \item Pericardiocentesis
    \item Thoracotomy
\end{enumerate}

\item The nurse is caring for a patient with a spinal cord injury at C4. Which complication should the nurse monitor for MOST closely?

\begin{enumerate}[label=\Alph*.]
    \item Autonomic dysreflexia
    \item Respiratory failure
    \item Deep vein thrombosis
    \item Pressure ulcers
\end{enumerate}

\item A patient with T4 spinal cord injury suddenly develops a blood pressure of 210/118 mmHg, severe headache, and facial flushing. The nurse's FIRST action should be:

\begin{enumerate}[label=\Alph*.]
    \item Administer prescribed antihypertensive
    \item Sit the patient upright
    \item Check for bladder distension
    \item Call the rapid response team
\end{enumerate}

\item \textbf{Select all that apply.} Which are components of the ``Trauma Triad of Death''?

\begin{enumerate}[label=\Alph*.]
    \item Hypothermia
    \item Hyperglycemia
    \item Acidosis
    \item Coagulopathy
    \item Hypoxia
\end{enumerate}

\item A patient arrives after an MVC with a seatbelt sign across the abdomen. The nurse should be MOST concerned about injury to which organs?

\begin{enumerate}[label=\Alph*.]
    \item Lungs and heart
    \item Small bowel and lumbar spine
    \item Kidneys and bladder
    \item Liver and spleen only
\end{enumerate}

\item The FAST exam is positive in all four quadrants. The nurse anticipates:

\begin{enumerate}[label=\Alph*.]
    \item CT scan of the abdomen
    \item Diagnostic peritoneal lavage
    \item Emergent laparotomy
    \item Serial abdominal examinations
\end{enumerate}

\item A patient with a Grade IV splenic laceration is being managed non-operatively. Which finding would indicate failure of conservative management?

\begin{enumerate}[label=\Alph*.]
    \item Stable hemoglobin over 24 hours
    \item Heart rate increasing from 88 to 124 bpm
    \item Decreasing abdominal pain
    \item Urine output of 45 mL/hour
\end{enumerate}

\item The nurse is caring for a patient with flail chest. Which intervention is MOST appropriate for initial management?

\begin{enumerate}[label=\Alph*.]
    \item Tape the flail segment to stabilize it
    \item Position patient with injured side up
    \item Administer adequate pain control and monitor for respiratory failure
    \item Immediately intubate the patient
\end{enumerate}

\item A patient has a pelvic fracture with suspected hemorrhage. Which intervention should the nurse anticipate FIRST?

\begin{enumerate}[label=\Alph*.]
    \item Application of a pelvic binder
    \item CT scan of the pelvis
    \item Foley catheter insertion
    \item External fixation in the OR
\end{enumerate}

\item \textbf{Select all that apply.} Which are contraindications to Foley catheter insertion in a trauma patient?

\begin{enumerate}[label=\Alph*.]
    \item Blood at the urethral meatus
    \item Scrotal hematoma
    \item High-riding prostate on rectal exam
    \item Pelvic fracture without other signs
    \item Perineal ecchymosis
\end{enumerate}

\item A trauma patient has a GCS of 6. The nurse knows that this patient:

\begin{enumerate}[label=\Alph*.]
    \item Can protect their own airway
    \item Requires close monitoring but no airway intervention
    \item Requires definitive airway management
    \item Should be observed for improvement before intubation
\end{enumerate}

\item The Cushing's triad (hypertension, bradycardia, irregular respirations) indicates:

\begin{enumerate}[label=\Alph*.]
    \item Spinal shock
    \item Increased intracranial pressure with impending herniation
    \item Cardiac tamponade
    \item Tension pneumothorax
\end{enumerate}

\item A patient with traumatic brain injury has an ICP of 28 mmHg and MAP of 85 mmHg. What is the cerebral perfusion pressure (CPP)?

\begin{enumerate}[label=\Alph*.]
    \item 113 mmHg
    \item 57 mmHg
    \item 28 mmHg
    \item 85 mmHg
\end{enumerate}

\item The target CPP for a patient with traumatic brain injury is:

\begin{enumerate}[label=\Alph*.]
    \item Greater than 40 mmHg
    \item Greater than 50 mmHg
    \item 60-70 mmHg
    \item Greater than 90 mmHg
\end{enumerate}

\item \textbf{Select all that apply.} Which interventions help reduce intracranial pressure?

\begin{enumerate}[label=\Alph*.]
    \item Elevate head of bed 30 degrees
    \item Keep head midline
    \item Administer hypertonic saline or mannitol
    \item Hyperventilate to PaCO2 of 25 mmHg routinely
    \item Prevent hyperthermia
    \item Maintain adequate sedation
\end{enumerate}

\item A burn patient has burns to the entire anterior trunk and both entire arms. Using the Rule of Nines, what is the estimated TBSA burned?

\begin{enumerate}[label=\Alph*.]
    \item 27\%
    \item 36\%
    \item 45\%
    \item 54\%
\end{enumerate}

\item A patient has a 40\% TBSA burn and weighs 80 kg. Using the Parkland formula, what is the total fluid requirement for the first 24 hours?

\begin{enumerate}[label=\Alph*.]
    \item 6,400 mL
    \item 9,600 mL
    \item 12,800 mL
    \item 16,000 mL
\end{enumerate}

\item Using the Parkland formula calculation from the previous question, how much fluid should be given in the first 8 hours?

\begin{enumerate}[label=\Alph*.]
    \item 3,200 mL
    \item 4,800 mL
    \item 6,400 mL
    \item 9,600 mL
\end{enumerate}

\item A patient with circumferential full-thickness burns to the chest is developing respiratory distress. The nurse anticipates which procedure?

\begin{enumerate}[label=\Alph*.]
    \item Intubation only
    \item Escharotomy
    \item Fasciotomy
    \item Skin grafting
\end{enumerate}

\item \textbf{Select all that apply.} Which findings suggest inhalation injury in a burn patient?

\begin{enumerate}[label=\Alph*.]
    \item Singed nasal hairs
    \item Carbonaceous sputum
    \item Facial burns
    \item Hoarse voice
    \item Burns occurred in enclosed space
    \item Elevated carboxyhemoglobin level
\end{enumerate}

\item The nurse is caring for a patient with electrical burns. Which complication should be monitored for MOST closely?

\begin{enumerate}[label=\Alph*.]
    \item Infection
    \item Cardiac arrhythmias
    \item Hypothermia
    \item Respiratory failure
\end{enumerate}

\item A patient presents after a house fire with cherry-red skin color and altered mental status. The nurse suspects:

\begin{enumerate}[label=\Alph*.]
    \item Cyanide poisoning
    \item Carbon monoxide poisoning
    \item Severe burn shock
    \item Anaphylaxis
\end{enumerate}

\item The treatment for carbon monoxide poisoning includes:

\begin{enumerate}[label=\Alph*.]
    \item Low-flow oxygen via nasal cannula
    \item 100\% oxygen via non-rebreather mask or hyperbaric oxygen
    \item Hydroxocobalamin administration
    \item Methylene blue administration
\end{enumerate}

\end{enumerate}

\newpage
\subsection*{Practice Exam Answer Key}

\begin{enumerate}
    \item \textbf{B} - Airway with cervical spine protection is always first in trauma
    \item \textbf{C, F} - Accessory muscle use and paradoxical movement indicate respiratory distress
    \item \textbf{C} - Tachycardia with weak, thready pulse indicates shock
    \item \textbf{C} - Eye 1 + Verbal 1 + Motor 4 (withdrawal) = 6
    \item \textbf{B} - Ipsilateral pupil dilation indicates CN III compression from herniation
    \item \textbf{D} - Both complete exposure and hypothermia prevention are essential
    \item \textbf{A, B, C, D, E} - All of these patients meet criteria for spinal immobilization
    \item \textbf{B} - Hypotension with bradycardia = neurogenic shock
    \item \textbf{B} - Tension pneumothorax is a clinical diagnosis; don't wait for X-ray
    \item \textbf{D} - Class IV: HR $>$140, severe hypotension, confused/lethargic
    \item \textbf{B} - 1:1:1 ratio for balanced resuscitation
    \item \textbf{A, B, C, D, E, F} - All are signs of successful decompression
    \item \textbf{C} - Beck's triad indicates cardiac tamponade; treatment is pericardiocentesis
    \item \textbf{B} - C4 injury affects diaphragm innervation (C3-5); respiratory failure is primary concern
    \item \textbf{B} - Sitting upright is FIRST action to use gravity to lower BP
    \item \textbf{A, C, D} - Hypothermia, acidosis, coagulopathy
    \item \textbf{B} - Seatbelt sign associated with small bowel and lumbar spine injuries
    \item \textbf{C} - Hemodynamically unstable patient with positive FAST needs emergent surgery
    \item \textbf{B} - Increasing heart rate indicates ongoing hemorrhage
    \item \textbf{C} - Pain control and monitoring; taping is contraindicated
    \item \textbf{A} - Pelvic binder reduces pelvic volume and tamponades bleeding
    \item \textbf{A, B, C, E} - All except pelvic fracture alone are contraindications
    \item \textbf{C} - GCS $\leq$8 requires definitive airway
    \item \textbf{B} - Cushing's triad indicates brainstem compression
    \item \textbf{B} - CPP = MAP - ICP = 85 - 28 = 57 mmHg
    \item \textbf{C} - Target CPP is 60-70 mmHg
    \item \textbf{A, B, C, E, F} - Routine hyperventilation is NOT recommended
    \item \textbf{B} - Anterior trunk (18\%) + both arms (9\% × 2 = 18\%) = 36\%
    \item \textbf{C} - 4 mL × 80 kg × 40\% = 12,800 mL
    \item \textbf{C} - Half in first 8 hours = 6,400 mL
    \item \textbf{B} - Escharotomy releases constriction from circumferential burns
    \item \textbf{A, B, C, D, E, F} - All are signs of inhalation injury
    \item \textbf{B} - Electrical burns cause cardiac arrhythmias; monitor on telemetry
    \item \textbf{B} - Cherry-red color and AMS suggest CO poisoning
    \item \textbf{B} - High-flow oxygen displaces CO from hemoglobin
\end{enumerate}

\newpage
\section*{Section 23: Final Summary and Key Takeaways}

\begin{tcolorbox}[colback=blue!5!white,colframe=blue!75!black,title=\textbf{TRAUMA NURSING: THE BIG PICTURE}]

\textbf{The Golden Hour Concept:}
\begin{itemize}
    \item Definitive care within 60 minutes of injury significantly improves survival
    \item Every minute counts in trauma resuscitation
    \item Systematic approach prevents missed injuries
\end{itemize}

\vspace{0.3cm}

\textbf{Primary Survey: ABCDE - Never Skip Steps}
\begin{itemize}
    \item \textbf{A}irway with C-spine protection - Secure airway first
    \item \textbf{B}reathing - Look, listen, feel; treat life threats immediately
    \item \textbf{C}irculation - Control hemorrhage; restore perfusion
    \item \textbf{D}isability - Neuro status; prevent secondary brain injury
    \item \textbf{E}xposure - See everything; prevent hypothermia
\end{itemize}

\vspace{0.3cm}

\textbf{Life-Threatening Conditions Requiring IMMEDIATE Intervention:}
\begin{itemize}
    \item Tension pneumothorax → Needle decompression
    \item Cardiac tamponade → Pericardiocentesis
    \item Massive hemothorax → Chest tube + OR
    \item Airway obstruction → Secure airway
    \item Uncontrolled hemorrhage → Direct pressure, tourniquet, OR
\end{itemize}

\vspace{0.3cm}

\textbf{Remember the Priorities:}
\begin{enumerate}
    \item Treat the greatest threat to life FIRST
    \item Don't get distracted by dramatic but non-life-threatening injuries
    \item Reassess constantly - patient condition changes rapidly
    \item Document thoroughly - legal and continuity of care
    \item Communicate clearly - SBAR, closed-loop communication
\end{enumerate}

\end{tcolorbox}

\vspace{0.5cm}

\begin{tcolorbox}[colback=red!5!white,colframe=red!75!black,title=\textbf{CRITICAL NUMBERS TO MEMORIZE}]

\begin{center}
\begin{tabular}{|l|l|}
\hline
\textbf{Parameter} & \textbf{Critical Value/Target} \\
\hline
GCS requiring intubation & $\leq$ 8 \\
\hline
Target CPP & 60-70 mmHg \\
\hline
ICP upper limit & $<$ 20-22 mmHg \\
\hline
Parkland formula & 4 mL × kg × \%TBSA \\
\hline
MTP ratio & 1:1:1 \\
\hline
Class III shock HR & 120-140 bpm \\
\hline
Class IV shock blood loss & $>$40\% ($>$2000 mL) \\
\hline
Permissive hypotension target & SBP 80-90 mmHg \\
\hline
Autonomic dysreflexia level & T6 and above \\
\hline
Target SpO2 in trauma & $\geq$94\% \\
\hline
\end{tabular}
\end{center}

\end{tcolorbox}

\vspace{0.5cm}

\begin{tcolorbox}[colback=green!5!white,colframe=green!75!black,title=\textbf{NURSING EXCELLENCE IN TRAUMA CARE}]

\textbf{The trauma nurse is:}
\begin{itemize}
    \item The patient's advocate when they cannot speak for themselves
    \item The eyes and ears detecting subtle changes
    \item The skilled hands performing life-saving interventions
    \item The calm presence in chaos
    \item The link between the patient and the trauma team
\end{itemize}

\vspace{0.3cm}

\textbf{Keys to Success:}
\begin{itemize}
    \item Know your protocols cold - there's no time to look things up
    \item Practice skills regularly - competence builds confidence
    \item Communicate clearly - lives depend on accurate information
    \item Take care of yourself - trauma nursing is demanding
    \item Debrief after difficult cases - learn and process
\end{itemize}

\vspace{0.3cm}

\textit{``In trauma, we don't have the luxury of time. We have the privilege of making every second count.''}

\end{tcolorbox}

\vspace{1cm}

\begin{center}
\textbf{\Large You are prepared. You are capable. You will save lives.}

\vspace{0.5cm}

\textit{Good luck on your exam!}
\end{center}

\end{document}